
\subsection{An abstraction for restricted \laVASS}

Our decision procedure for emptiness of restricted \laVASS is based on
reducing this problem to a reachability problem in a
carefully designed finite-state abstraction of the state-space of
\laVASS. Throughout this section, let $V=\langle Q,2d,
\Sigma,\Delta,\lambda,q_0,F,\Phi\rangle$ be a restricted \laVASS.  We
first define the state space of the abstracted \laVASS.
%
\begin{definition}
  An abstract configuration is a tuple
%
  \begin{multline*}
  \alpha=(q,
  m_1,n_1,\dots,m_{d},n_{d},u_1,u_2,\dots,u_{d-1},\ell_1,\dots,\ell_d)\\
  \in Q\times (\N\cup\{\bot\})^{2d} \times \N^{d-1}\times(\N \cup \{ \top \})^{d},
  \end{multline*}
  %
  such that $m_i,n_i \in [0,2dM_i] \cup \{ \bot \}$
  and $u_i \in [0,U_i]$ and $\ell_i\in[0,M_i-1] \cup \{\top\}$,
  where
  \begin{itemize}
	\item $M_i \defeq \floor{\abs Q^{((1/8)\cdot 32^{i-1} + 1)}}$; and
	\item $U_i \defeq \abs Q^{(32^{i-1}+4)}$.
  \end{itemize}
\end{definition}
%
The idea is that $m_i,n_i$ store the residue classes modulo $\ell_i$
of the counter pair $x_i,y_i$ respectively, where the value $\top$ for
$\ell_i$ acts as an indicator that we are storing actual values and
not residue classes. The value $\bot$ for some $x_i$ or $y_i$
indicates that the counter has not yet been initialised. If for an
update function $f$, $f=\wPP$ or $f=\wTTPO$ then $f(\bot)\defeq
1$; otherwise $f(\bot) \defeq \bot$, and we stipulate that $\bot \bmod
n = \bot$. The value of $u_i$ in an abstract configuration carries the
current difference between the value of the counters $y_{i}$ and
$y_{i+1}$. This difference is potentially unbounded, however for our
purposes it suffices to only store its value if it is less than $U_i$,
and to indicate the fact that it is at least $U_i$ by the value
$u_i=U_i$.
%This difference must not become negative since $u_1,\dots,u_{d-1}$ are
%all natural numbers.

We denote the (finite) set of all abstract configurations of $V$ by
$A(V)$. Let us now define a transition relation
${\xrightarrow{\cdot}}\subseteq A(V) \times \Delta \times A(V)$ such
that $\alpha\xrightarrow{t}\alpha'$, $t=(q,a,q')\in \Delta$ if and
only if:
%
\begin{itemize}
    \item $\alpha=(q,m_1,n_1,\dots,u_{d-1},\ell_1,\dots,\ell_d)$ and $\alpha'=(q',m_1',n_1',\dots,u_{d-1}',\ell_1,\dots,\ell_d)$;
    \item $\lambda(t)=(f_{x_1},f_{y_1},\dots,f_{x_d},f_{y_d})$;
    \item $f_{x_i} = \wPP$ for all $i$ such that $m_i\neq \bot$;
    \item if $\ell_i\neq\top$, $m_i' = f_{x_i}(m_i)\bmod \ell_i$ and $n_i'=f_{y_i}(n_i)\bmod \ell_i$;
    \item if $\ell_i=\top$, $m_i' = f_{x_i}(m_i)$ and $n_i'=f_{y_i}(n_i)$; and
    \item for all $i\in\{1,\dots,d-1\}$,
      $$u_i' =
    \begin{cases*}
      min(2u_i+1,U_i) & if $f_{y_{i}}=\wTTPO, f_{y_{i+1}}=\wTT$ \\
      min(2u_i,U_i) & if $f_{y_{i}}=f_{y_{i+1}}$ \\
      min(2u_i-1,U_i) & if $f_{y_{i}}=\wTT, f_{y_{i+1}}=\wTTPO$. \\
    \end{cases*}
    $$
\end{itemize}
%
Assuming that the value of $y_i$ is at least the value of $y_{i+1}$,
which we will always ensure, the definition of how to update $u_i$
ensures that it exactly stores the difference $y_i-y_{i+1}$ unless the
difference becomes too large, in which case it is levelled off at
$U_i$.

An abstract configuration path is a sequence of abstract
configurations and transitions of the form
$R=\alpha_1\xrightarrow{t_1} \alpha_2 \xrightarrow{t_2} \cdots
\xrightarrow{t_{n-1}} \alpha_n$.

Given two consecutive $y$-counters $y_i,y_{i+1}$ and $\delta_i\in \N$,
we say that $y_{i}-y_{i+1}=\delta_i$ and $y_{i}-y_{i+1}\geq \delta_i$
are \emph{$y$-constraints}. Let $Y$ be a set of $y$-constraints, an
abstract configuration
$\alpha=(q,m_1,n_1,\dots,u_1,\dots,u_{d-1},\ell_1,\dots,\ell_d)$
\emph{respects $Y$} whenever
%
\begin{itemize}
\item
  $u_i\geq \delta_i$
  for all constraints of type $y_{i}-y_{i+1} \geq \delta_i$ in $Y$,
\item and $u_i<U_i$ and $u_i=\delta_i$ for all constraints
  $y_{i}-y_{i+1} = \delta_i$ in $Y$.
\end{itemize}
%
We say that $\alpha_f$ is a \emph{final abstract configuration
  respecting $Y$} whenever $q\in F$, $m_i=n_i$ for all $1\leq i\leq
d$, and $\alpha_f$ respects $Y$.


\subsection{Witnessing certificates}

While any concrete accepting run of an \laVASS gives rise to an
abstract configuration path ending in an accepting abstract
configuration, the converse does not hold. This motivates the
introduction of reachability and witnessing certificates, which are
special abstract configuration paths that carry further information
that eventually enables us to derive from a witnessing certificate a
concrete accepting run of an \laVASS.

A \emph{reachability certificate} is a tuple $(R,X,Y,L)$ such that
$R=\alpha_1\xrightarrow{t_1} \alpha_2 \xrightarrow{t_2} \cdots
\xrightarrow{t_{n-1}} \alpha_n$ is an abstract configuration path, and
$X,Y,L\colon \{1,\ldots,d\} \to \{1,\ldots,n\}$. Here, $X(i)$ and
$Y(i)$ indicate the position where the $x_i$-counter and $y_i$-counter
obtain a value different from $\bot$ for the first time. Moreover,
$L(i)$ is the position where a loop of length $\ell_i$ can be found.
Formally, $(R,X,Y,L)$ is required to have the following properties:
%
\begin{enumerate}[(a)]
\item $\alpha_1 = (q_0,\bot, \ldots, \bot,0,\ldots,0,
  \ell_1,\ldots,\ell_d)$ and if $\ell_i=\top$ then $\ell_j=\top$ for
  all $i<j\le d$;
\item $\lambda(x_i,t_{X(i)-1})=\wPP$ and $\lambda(y_i,t_{Y(i)-1})=\wTTPO$ for all $1\le i\le d$;
\item $\lambda(x_i,t_{j})=\wID$ for all $1\le j < X(i)-1$;
\item $\lambda(y_i,t_{j})=\wTT$ for all $1\le j < Y(i)-1$;
\item $X,Y,L$ are monotonic;
\item for all $1\le i\le d$, if $\ell_i\neq \top$ then
  \begin{itemize}
  \item $X(i)\le L(i)< Y(i)$; and
  \item there is a simple $\alpha_{L(i)}$-loop $\alpha_{L(i)}
    \xrightarrow{t_1'} \alpha'_2 \xrightarrow{t_2'} \cdots
    \xrightarrow{t_{\ell_i-1}} \alpha'_{\ell_i-1}
    \xrightarrow{t_{\ell_i}'} \alpha_{L(i)}$ of length $\ell_i$.
  \end{itemize}
\end{enumerate}
%
Those conditions can be interpreted as follows. Condition~(a) asserts
that the certificate starts in an initial abstract configuration. We
require that $\top$ monotonically propagates since since the absence of
a loop for counter $x_i$ implies that the remainder of a path is
short, hence we can afford to subsequently store actual counter
values and not residue classes. Conditions~(b), (c) and~(d) assert
that $X(i)$ and $Y(i)$ are the first position where the counters
$x_i,y_i$ hold value different from $\bot$. Condition~(e) states that
the counters $x_{i+1}$, $y_{i+1}$ do not carry a value different from
$\bot$ before the counters $x_{i}$ and $y_{i}$,
respectively. Condition~(f) implies that, if $\ell_i\neq \top$ then
between the first update for counter $x_i$ and the first update for
counter $y_i$ there is a position $L(i)$ where we can find a loop in
the abstract configurations of length $\ell_i$. Notice that if
$x_j=\bot$ or $y_j=\bot$ in $\alpha_{L(i)}$ then $x_j$ and $y_j$
remain to hold $\bot$ along this loop, i.e., this loop does not update
counters that have not been initialised already.

Given $R$, the set of $y$-constraints induced by $R$ is the smallest
set containing
%
\begin{itemize}
\item $y_{i}-y_{i+1}\geq \delta_i$, where $\delta_i \defeq
  X(i+1)-X(i)$ if there is a $j$ such that $X(i)\le L(j) < X(i+1)$;
  and
\item otherwise $y_{i}-y_{i+1} = \delta_i$, where $\delta_i \defeq
  X(i+1)-X(i)$,
\end{itemize}
for all $1\le i< d$ such that $\ell_i\neq\top$.

We introduce some further notation. Given a reachability certificate
$R$, we denote by $\pi(R)$ the run corresponding to $R$ in the
configuration graph of $V$, with the initial configuration
$(q_0,0,0,\ldots,0,0)$. Given indices $1\le i\le j\le n$, we denote by $R[i,j]$ the
segment $\alpha_i \xrightarrow{t_i} \alpha_{i+1} \cdots
\xrightarrow{t_{j-1}} \alpha_j$ of $R$, and by $R[i]\defeq \alpha_i$.
We say that $R$ is a \emph{witnessing certificate} if, for $a\le d$
being the largest index such that $\ell_a \neq \top$:
%
\begin{itemize}
\item $R[1,Y(a)]$ is a simple path and $n-Y(a)\le 2dM_{d+1}$;
\item $\alpha_n$ is a final abstract configuration respecting the set
  of induced $y$-constraints; and
\item $\val(\pi(R),x_a) \le \val(\pi(R),y_a)$.
\end{itemize}
%
Sometimes we will speak of witnessing certificates \emph{restricted}
to a set of counters. By that we mean a witnessing certificates where
the relevant Conditions~(a)--(f) are only required for that set of
counters.

Now we are ready to provide a proof for \Cref{prop:main}, that stated that language emptiness for restricted \laVASS can be decided in $\mathrm{NSPACE}( \abs V \cdot 2^{O(d)})$.
\begin{proof}[Proof of \Cref{prop:main}]
  Clearly, an abstract configuration can be stored in space $\abs V \cdot 2^{O(d)}$. An NEXPSPACE algorithm can hence
  non-deterministically choose an initial configuration and
  non-deterministically verify that it leads to a final abstract
  configuration along a path that is a witnessing certificate. To this
  end, the algorithm computes the set of induced $y$-constraints
  on-the-fly while guessing the reachability certificate, and verifies
  them in the last configuration. Note that the $y$-constraints can be stored in space $\abs V \cdot 2^{O(d)}$. Finally, the requirement
  $\val(\pi(R),x_a) \le \val(\pi(R),y_a)$ can also be verified in
  exponential space since we require that $R[1,Y(a)]$ is a simple path and $n-Y(a)\le 2M_{d+1}$.
\end{proof}
%
In the next section we argue the correctness of our algorithm by proving the following theorem:
\begin{theorem}\label{theorem:certificate iff non-empty language}
The language of a restricted \laVASS $V$ is non-empty if and only if there exists a witnessing certificate for $V$.
\end{theorem}
%
